% !TEX root = ../mainfile.tex
% !TEX encoding = UTF-8 Unicode
% !TEX spellcheck = en_US

\chapter{Introduction}
\label{ch:introduction}

Within a single year, two major shocks moved global markets in opposite directions for the
U.S.\ dollar. On 2~April 2025---``Liberation Day''---the United States announced sweeping
tariffs on virtually all of its trading partners, and the dollar \emph{fell}: it depreciated
against the euro, the yen and the Swiss franc even as equities sold off and volatility
spiked, the textbook conditions under which the dollar normally \emph{rises}. Less than a
year later, as escalation in the Gulf threatened to close the Strait of Hormuz and sent
crude sharply higher, the dollar did exactly what the safe-haven literature predicts---it
appreciated. Two episodes, both with the United States implicated, produced a mirror image.
This thesis asks why.

The puzzle speaks to a live debate about whether the dollar's exceptional status is eroding.
A growing literature argues that the currency's safe-haven premium has weakened, pointing to
the unusual co-movement of a falling dollar, rising Treasury yields and rising risk premia
during the 2025 tariff episode \parencite{corsetti2025,kamin2025,obstfeld2023,jiang2026}.
Against this, the long tradition on the dollar's ``exorbitant privilege'' and its role as the
world's reserve and safe-haven currency \parencite{gourinchasrey2005,habib2012,ranaldo2010}
implies the status should be durable. The two 2025--2026 shocks offer an unusually clean
natural experiment for adjudicating the debate, because they differ in \emph{kind}: one is a
U.S.\ trade-policy shock, the other a geopolitical oil-supply shock.

\section*{Research question}

The central question of this thesis is therefore: \emph{does the dollar's safe-haven
response depend on the type of shock---even when the United States is, in some sense, the
originator of both?} The design contrasts three episodes. The Liberation Day tariff
announcement is a purely U.S.-originated economic-policy shock. The 2026 Israel--Iran/Hormuz
crisis is a geopolitical oil-supply shock in which the United States is implicated as a
co-belligerent but still functions as a financial safe harbour. And the 2022 Russian
invasion of Ukraine---an external military/oil shock not authored by the United
States---serves as a control for the ``classical'' safe-haven response. Comparing the three
isolates whether the dollar's behaviour is conditioned by the nature of the shock or merely
by who caused it.

\section*{Approach and contribution}

The thesis answers the question with two complementary empirical layers on a single daily
dataset spanning January~2019 to June~2026. An \emph{event study} measures the cumulative
abnormal response of the dollar, a cross-section of currency baskets, gold, equities and
crude oil around each shock, and tests formally whether the dollar's response differs across
events. A complementary \emph{regression} then tests whether the dollar's relationship with
risk itself changed across the episodes---whether its safe-haven sensitivity to volatility
held or broke. The two layers reinforce one another: the event study shows that discrete,
dated shocks move the dollar in opposite directions depending on shock type, and the
regression shows that the dollar's safe-haven reflex broke down for the trade-policy shock
yet held for the geopolitical ones.

The contribution is threefold. First, to the best of the author's knowledge this is the
first systematic side-by-side comparison of the Liberation Day and Iran/Hormuz episodes for
currency markets. Second, the thesis bridges the safe-haven and commodity-currency
literatures by carrying the oil channel and a gold benchmark through the same design,
showing that gold behaves as a hedge \emph{against} the dollar rather than as an
unconditional substitute for it. Third, the rebuilt carry/dollar risk factors and the WMR
currency panel extend standard cross-sectional tools through to mid-2026, allowing the most
recent shocks to be studied with consistent, citable data.

The remainder of the thesis is organised as follows. Chapter~\ref{ch:literature} reviews the
safe-haven, dollar-erosion and commodity-currency literatures and states the predictions
that follow from them. Chapter~\ref{ch:background} narrates the 2025--2026 events.
Chapter~\ref{ch:methodology} describes the data and the event-study and regression
methodologies. Chapter~\ref{ch:event_study} presents the event-study results and
Chapter~\ref{ch:regression} the regression test. Chapter~\ref{ch:discussion} discusses the
findings and their implications,
and Chapter~\ref{ch:conclusion} concludes.
